\begin{figure}[!htbp]
\centering
\resizebox{\linewidth}{!}{%
\begin{tikzpicture}
\begin{axis}[
    xbar,
    width=130mm, height=85mm,
    scale only axis,
    xmin=0, xmax=100,
    enlarge x limits={upper, value=0.06},
    xlabel={Participants answering \emph{Yes} (\%)},
    symbolic y coords={
        Calendar,App info/perf.,Audio files,Health/fitness,App activity,
        Device IDs,Web browsing,Files/docs,Contacts,
        Photos/videos,Financial info,Messages,Location,Personal info},
    ytick=data, y dir=normal,
    xtick={0,20,40,60,80,100},
    nodes near coords, point meta=x,
    nodes near coords={\pgfmathprintnumber[fixed,precision=1]{\pgfplotspointmeta}\,\%},
    nodes near coords align=horizontal,
    every node near coord/.append style={font=\footnotesize, color=ngSlate,
                                         /pgf/number format/.cd, fixed, precision=1},
    y tick label style={font=\footnotesize},
    x tick label style={font=\footnotesize},
    bar width=6.5pt, enlarge y limits=0.04,
    axis line style={ngSlate}, tick style={ngSlate},
    grid=major, grid style={ngSlate!18}
]
\addplot+[draw=ngBlue, fill=ngBlue!85, line width=0.4pt] coordinates {
    (78.2,Personal info) (75.9,Location)   (74.7,Messages)
    (69.0,Financial info)(69.0,Photos/videos)(63.2,Contacts)
    (62.1,Files/docs)    (60.9,Web browsing)(60.9,Device IDs)
    (55.2,App activity)  (51.7,Health/fitness)(50.6,Audio files)
    (46.0,App info/perf.)(42.5,Calendar)
};
\end{axis}
\end{tikzpicture}%
}
\caption{Perceived personalness of Google Play Data Safety categories
(\textit{N}\,=\,87). Categories range from broadly recognised personal
data (top) to deeply ambiguous (bottom). The 35-point spread is itself
the contribution: a flat category list hides this gradient from users.}
\label{fig:rq1-personalness}
\end{figure}
